\documentclass[multi=tikzpicture,border=8pt]{standalone}
\usepackage[UTF8]{ctex}
\usepackage{amsmath,amssymb,tikz}
\usetikzlibrary{arrows.meta,calc}
\definecolor{blue}{HTML}{17365D}
\definecolor{red}{HTML}{AC3931}
\tikzset{every node/.style={font=\small},arrow/.style={-{Stealth[length=2.8mm]},thick},dim/.style={|<->|,thin},box/.style={draw,rounded corners,align=center,minimum width=2.5cm,minimum height=.8cm}}
\begin{document}
\begin{tikzpicture}

\draw[arrow](0,-.4)--(0,5)node[above]{实空间轴$z$};\draw[blue,very thick,domain=0:720,samples=160]plot({.65*cos(\x)},{\x/180});\draw[red,very thick,domain=0:720,samples=160]plot({.65*cos(\x+135)},{\x/180});\draw[dim](1.1,0)--(1.1,2)node[right,midway]{$h$};\draw[dim](-.65,-.8)--(.65,-.8)node[below,midway]{$2r$};\begin{scope}[shift={(5,2)}]\draw[arrow](-1.9,0)--(1.9,0)node[right]{$q_\perp$};\draw[arrow](0,-2)--(0,2.6)node[above]{$q_z$};\draw[blue,very thick](-1.5,-1.788)--(1.5,1.788);\draw[blue,very thick](-1.5,1.788)--(1.5,-1.788);\draw(.55,.656)arc(50:130:.856);\node[above]at(0,1){$80^\circ$};\end{scope}

\end{tikzpicture}
\begin{tikzpicture}

\draw[very thick,blue](-2,-1.7)--(2,1.7)node[above right]{局部切线};\draw[arrow,red](0,0)--(-1.5,1.765)node[left]{衍射法线};\draw[arrow](-2.5,0)--(2.7,0)node[right]{横向};\draw[arrow](0,-2.2)--(0,2.5)node[above]{轴向};\node[align=center]at(4.4,-1.3){切线斜率$h/(2\pi r)$\\衍射臂斜率$2\pi r/h$};\draw[thick](-.23,.28)--(.05,.52)--(.28,.24);

\end{tikzpicture}
\begin{tikzpicture}

\draw[arrow](-2.3,0)--(2.4,0)node[right]{$q_\perp$};\draw[arrow](0,-2.6)--(0,2.7)node[above]{$q_z$};\foreach \j in {-3,-2,-1,0,1,2,3}{\draw[blue](-1.8,\j*.7)--(1.8,\j*.7);}\draw[dim](2.1,.7)--(2.1,1.4)node[right,midway]{$2\pi/h$};\node[below]at(0,-3.1){$h\sin\theta\simeq h\theta=\lambda$};

\end{tikzpicture}
\begin{tikzpicture}

\draw[gray](0,0)circle(1.5);\draw[arrow,blue](0,0)--(1.5,0)node[right]{$A_n$};\draw[arrow,red](0,0)--(-1.5,0)node[left]{$A_ne^{in\phi}$};\node[above]at(0,2){第四级相消：$4\phi=(2j+1)\pi$};\node[below]at(0,-2){$\phi=45^\circ,135^\circ,225^\circ,315^\circ$};

\end{tikzpicture}
\begin{tikzpicture}
\draw[arrow](0,0)--(8.3,0)node[right]{$n$};\draw[arrow](0,0)--(0,3.3)node[above]{$I_n/(4|A_n|^2)$};\fill[blue](1,2.13388)circle(2pt);\fill[red](1,0.36612)circle(2pt);\fill[blue](2,1.25000)circle(2pt);\fill[red](2,1.25000)circle(2pt);\fill[blue](3,0.36612)circle(2pt);\fill[red](3,2.13388)circle(2pt);\fill[blue](4,0.00000)circle(2pt);\fill[red](4,0.00000)circle(2pt);\fill[blue](5,0.36612)circle(2pt);\fill[red](5,2.13388)circle(2pt);\fill[blue](6,1.25000)circle(2pt);\fill[red](6,1.25000)circle(2pt);\fill[blue](7,2.13388)circle(2pt);\fill[red](7,0.36612)circle(2pt);\fill[blue](8,2.50000)circle(2pt);\fill[red](8,2.50000)circle(2pt);\node[left]at(0,2.5){1};\draw[dashed](4,0)--(4,2.8);\node[below]at(4,0){4};\node[above]at(4,3.3){两分支都在第四级相消};\node[blue]at(2.3,3){$\phi=45^\circ$};\node[red]at(6.7,3){$\phi=135^\circ$};
\end{tikzpicture}
\end{document}
